\documentclass[11pt,a4paper]{article}
% ---------------------------------------------------------------------------------------------
%  First user guide of the ArgonCrModel echelle pipeline (Main.py + pipeline/).
%  Plain LaTeX, standard packages only: compiles with pdflatex, e.g. on Overleaf.
% ---------------------------------------------------------------------------------------------
\usepackage[T1]{fontenc}
\usepackage[utf8]{inputenc}
\usepackage{lmodern}
\usepackage[margin=2.3cm]{geometry}
\usepackage{booktabs}
\usepackage{tabularx}
\usepackage{array}
\usepackage{enumitem}
\usepackage{xcolor}
\usepackage{listings}
\usepackage[hidelinks]{hyperref}

\setlist{itemsep=2pt,topsep=4pt}
\setlength{\parskip}{4pt}
\setlength{\parindent}{0pt}
\renewcommand{\arraystretch}{1.25}
\newcolumntype{L}{>{\raggedright\arraybackslash}X}

% inline code / file names (underscores and braces are printed as they are; do not put % or # inside)
\newcommand{\code}[1]{\texttt{\detokenize{#1}}}

\definecolor{codebg}{gray}{0.95}
\definecolor{comment}{rgb}{0.35,0.45,0.35}
\lstset{basicstyle=\ttfamily\small, backgroundcolor=\color{codebg}, frame=single, rulecolor=\color{gray!50},
        columns=fullflexible, keepspaces=true, breaklines=true, showstringspaces=false,
        commentstyle=\color{comment}, xleftmargin=3pt, xrightmargin=3pt, aboveskip=6pt, belowskip=6pt}
\lstdefinelanguage{toml}{morecomment=[l]{\#}, morestring=[b]"}
\lstdefinelanguage{shell}{morecomment=[l]{\#}}

\newenvironment{note}{\par\smallskip\noindent\begin{minipage}{\linewidth}\small\textbf{Note.}~}{\end{minipage}\par\smallskip}

\title{ArgonCrModel: a first user guide\\[2pt]\large CR-model fits of echelle spectra of Ar and Ar/admixture microwave plasmas}
\date{October 2026}

\begin{document}
\emergencystretch=3em                 % long file names: stretch lines rather than run into the margin
\maketitle
\tableofcontents
\bigskip

% =============================================================================================
\section{What the code does}
% =============================================================================================
You give the code \textbf{echelle emission spectra} of an argon microwave plasma (2.42\,GHz), with or without a
second gas such as N$_2$, and the \textbf{laser-absorption (LAS) table} of the same discharge. For every plasma
condition (pressure, power, admixture \%) it does the following:
\begin{enumerate}
  \item fits the emission lines of every spectrum (line areas);
  \item takes the gas temperature $T_g$ of the condition from the LAS measurements, and builds, if it does not
        exist yet, the BOLSIG+ EEDF library and the CR-model table that the condition needs;
  \item fits the Ar\,I line intensities with the collisional--radiative (CR) model. This gives the E/N of the EEDF,
        the effective electron temperature $T_{e,\mathrm{eff}}$, the electron density $N_e$ (fitted or pinned) and
        the densities of all Ar levels, including the 1s metastables, which are compared with LAS;
  \item for conditions with N$_2$: does N\,I/Ar\,I actinometry, giving the N-atom fraction $n_N/n_{Ar}$, the
        N-atom density and the N$_2$ dissociation degree;
  \item writes tables and figures, and returns everything as one \textbf{results dict} (\S\ref{sec:results}).
\end{enumerate}
Everything is controlled by \textbf{one input file} (TOML). You do not need to edit any Python to run the model.

% =============================================================================================
\section{Setting up}
% =============================================================================================
\textbf{Python} 3.11 or newer (the input file is read with \code{tomllib}), with \code{numpy}, \code{scipy},
\code{pandas} and \code{matplotlib}. Spyder works well but is not required. Tested with Python 3.11.9, numpy 2.2,
scipy 1.15, pandas 2.3 and matplotlib 3.10.

\textbf{BOLSIG+} (\code{Bolsig/bolsigminus.exe}, Windows) is used only when a new EEDF library has to be built.
Already-built libraries are stored in \code{InputData/Bolsig/} and are reused.

\textbf{Paths.} Every path is relative to the folder that holds \code{Main.py} (the repo folder). The only
exception is a full path you write in the input file, for example for the echelle data.

\begin{table}[h]
\centering\small
\begin{tabularx}{\linewidth}{@{}lL@{}}
\toprule
Folder / file & What it is \\
\midrule
\code{Main.py} & \textbf{the script you run} \\
\code{input_example.toml} & example input file: \textbf{every setting with its default and explanation} \\
\code{InputFiles/} & your own input files (suggested place) \\
\code{CreateEEDFScheme.py} & tool to define new EEDF options (\S\ref{sec:eedf}) \\
\code{pipeline/} & the code of the pipeline; \code{pipeline/README.md} explains each file \\
\code{Experimental_Data/EchelleData/} & echelle spectra, one folder per data set \\
\code{Experimental_Data/LAS_Data/} & the LAS table(s) \\
\code{Experimental_Data/Output/} & results, and the caches (CR tables, logs) \\
\code{InputData/} & atomic data, cross sections, quenching tables; \code{Bolsig/} = EEDF libraries \\
\code{Bolsig/} & BOLSIG+ program and its cross-section files \\
\code{SmallAnalysisScripts/}, \code{Scripts/} & stand-alone analysis and test scripts (not needed for a run) \\
\code{LatexReport/} & the LaTeX report the pipeline can refresh (\code{[report]}) \\
\bottomrule
\end{tabularx}
\end{table}

% =============================================================================================
\section{Quick start}
% =============================================================================================
\begin{enumerate}
  \item Copy \code{input_example.toml} to, for example, \code{InputFiles/my_run.toml}.
  \item In \code{[run]}, set \code{output_folder} to a \textbf{new} folder, so you do not overwrite earlier results:
\begin{lstlisting}[language=toml]
[run]
output_folder = "Experimental_Data/Output/MyFirstRun"
\end{lstlisting}
  \item Describe your data sets in \code{[[dataset]]} blocks (\S\ref{sec:datasets}). Switch off the example data
        sets you do not need with \code{run = false}.
  \item \textbf{Always check first} with \code{--plan}. This reads and checks the input file, lists the conditions
        with their $T_g$ and LAS 1s5, and tells you which libraries and CR tables would be built. Nothing is
        fitted or overwritten:
\begin{lstlisting}[language=shell]
python Main.py InputFiles/my_run.toml --plan
\end{lstlisting}
  \item Run it:
\begin{lstlisting}[language=shell]
python Main.py InputFiles/my_run.toml
\end{lstlisting}
  \item Look at \code{<output_folder>/fits_all.csv}, the figures in \code{<output_folder>/<data set>/figures/},
        and the results dict \code{RES} (\S\ref{sec:results}).
\end{enumerate}

\textbf{In Spyder:} open \code{Main.py}, then Run $\rightarrow$ Configuration per file $\rightarrow$
\emph{Command line options}: \code{InputFiles/my_run.toml} (add \code{--plan} to check). After the run,
\code{RES} is in the Variable Explorer. From the IPython console you can also call
\code{import Main; RES = Main.main(["InputFiles/my_run.toml"])}.

% =============================================================================================
\section{Your input data}
% =============================================================================================
\subsection{Echelle spectra}
\begin{itemize}
  \item One folder per data set (one sweep), normally inside \code{Experimental_Data/EchelleData/}. It contains
        the \code{.spa} exports of the echelle software, several repeats per condition.
  \item \textbf{Intensity calibration:} the third header flag of a \code{.spa} file says whether it was exported
        with the intensity calibration. The pattern \code{[spectra] export_flags} (\code{?} = anything) decides
        which spectra are used. \code{"?-?-1-?"} or \code{"?-?-1-1"} keeps only calibrated spectra. Uncalibrated
        spectra are listed in the printout and skipped.
  \item Files that match no file-name pattern, such as backgrounds (\code{ref1.spa}), are ignored.
  \item A condition needs at least \code{[spectra] min_repeats} spectra (default 2).
\end{itemize}

\subsection{Describing a data set: the \texttt{[[dataset]]} block}\label{sec:datasets}
Each data set is one block. There are two ways to tell the code which file belongs to which condition.

\textbf{(a) A file-name pattern} with \code{{x}} = the swept value and \code{{rep}} = the repeat number. The
decimal point of \code{{x}} can be written \code{_} or \code{.}. For example, \code{Percent_0_33__2.spa} is
read as $x = 0.33$, repeat 2.
\begin{lstlisting}[language=toml]
[[dataset]]
name = "80W_1Torr"                    # short name: output folder and legend
folder = "N21TorrAdmixtureSweep80W"   # relative to [run] echelle_folder, or a full path
files = "Percent_{x}__{rep}"          # file name without .spa; several: ["pattern1", "pattern2"]
sweep = "N2"                          # what {x} is: "N2" (% in the feed), "pressure" (mTorr), "power" (W)
power_W = 80                          # give the TWO conditions that are not swept
pressure_mTorr = 1000
\end{lstlisting}

\textbf{(b) A file list} (CSV): use this when the file names have no clean pattern. Write
\code{file_list = "my_files.csv"}, with the path relative to the input file. Then \code{power_W},
\code{pressure_mTorr} and \code{N2_percent} are not needed in the block. The example is
\code{input_files_80W_example.csv}:
\begin{lstlisting}
file,pressure_mTorr,gas1,gas2,gas2_percent,power_W,flow_sccm,repeat,note
Percent_0_33__1,1000,Ar,N2,0.33,80,,,
Percent_0_33__2,1000,Ar,N2,0.33,80,,,
\end{lstlisting}
Required columns: \code{file}, \code{pressure_mTorr}, \code{gas1}, \code{gas2}, \code{gas2_percent},
\code{power_W}. Optional: \code{flow_sccm}, \code{repeat}, \code{Tg_K}, \code{note}. Lines starting with the
hash sign are comments.

\textbf{Optional keys} of a block: \code{run = false} skips the data set (it stays in the file);
\code{exclude_x = [0.0, 1.5]} leaves out these sweep values; \code{Tg_K = 600} uses a fixed gas temperature
instead of the LAS trend. Pure-Ar data sets (0\,\% admixture) get the CR fit and the metastables but no
actinometry.

\subsection{The LAS table}
\code{[las] file} (default \code{Experimental_Data/LAS_Data/Master_AllMeasurements.csv}) is used for three things:
the \textbf{gas temperature} of every condition (a trend $T_g(p, P, \%)$ fitted to all LAS rows, rounded to
\code{[eedf] Tg_step_K}), the \textbf{LAS 1s5 density} shown against the CR model, and an optional
\textbf{summary table} with the OES results added. Both layouts are read: one row per measurement
(\code{Master_AllMeasurements}) or one row per condition (\code{Master_Summary}). The columns used are
\code{run_id}, \code{date}, \code{pressure} [Torr], \code{power} [W], \code{Nitrogen_Percent}, \code{frequency}
[MHz], \code{T_gas_K} (and \code{_err}), and \code{N_s_cm3} (the 1s5 density, and \code{_err}). Useful settings
are \code{frequency_MHz}, \code{match_power_W}, \code{runs = "latest"} or \code{"all"}, and
\code{exclude_runs = ["<run id>"]}.

\subsection{Data you normally do not touch}
These live in \code{InputData/}: the Ar level and transition lists (\code{ArgonLevelList.json},
\code{ArgonReactionList.json}), the CR cross sections (default \code{ArgonCrossSections_BSR_analytic5p.json}),
the radiation-trapping table, and the two quenching tables \code{Ar_1s_quenching_data.csv} (4s levels) and
\code{Ar_2p_quenching_Sadeghi2001.csv} (4p levels). To add quenching rates for a new gas, add rows to these CSV
files.

% =============================================================================================
\section{The input file}
% =============================================================================================
\code{input_example.toml} lists \textbf{every} setting with its default and a full explanation. Read it once.
A setting you leave out keeps its default. A misspelled setting or section is an error, and the message lists
the allowed names. TOML basics: \texttt{\#} starts a comment, text goes in "double quotes", \code{true}/\code{false}
are lower case, lists are written \code{[1, 2, 3]}.

\begin{table}[h]
\centering\small
\begin{tabularx}{\linewidth}{@{}lL@{}}
\toprule
Section & Main settings (default) \\
\midrule
\code{[run]} & \code{output_folder}; \code{echelle_folder}; \code{overwrite} (true: clears old results of a data set, line fits kept); \code{show_plots} (false); \code{plan_only} \\
\code{[model]} & CR-model physics: \code{cross_sections}, \code{plasma_radius_m} (0.04), \code{trap_lines} (\code{"all"}), \code{escape_mode} (\code{"table"}), \code{atom_transfer}, \code{quench_4p}, \code{N2_quench_scale} (1), \code{N2_dissociation} (0), \code{admixture_gas} (\code{"N2"}) \\
\code{[eedf]} & \code{model} (\code{"pure_Ar"}, \code{"ArN2_vib_saturated"}, \code{"ArN2_vib_Tv"} or an EEDF scheme name); \code{Tv_K}, \code{Tv_grid_K}; \code{Tg_step_K} (50) \\
\code{[fit]} & \code{Ne_pin} (\code{"free"}, \code{"free_else_nc"}, \code{"nc"} or a density); \code{nc_penalty}; \code{sigma_model} (0.20); line selection: \code{wavelength_ranges_nm}, \code{exclude_lines_nm}, \code{min_rating}, \code{max_rel_err}, \code{merge_nm} \\
\code{[spectra]} & \code{export_flags}; \code{min_repeats} (2); \code{line_status} \\
\code{[actinometry]} & extra N\,I lines, excluded Ar references, outlier fences, cross-section errors \\
\code{[plots]} & one true/false per figure type; \code{dpi} \\
\code{[las]} & the LAS table and how it is matched to the OES conditions; the summary table \\
\code{[report]} & \code{make = true} refreshes \code{LatexReport/figures} and \code{tables} after the run \\
\code{[[dataset]]} & one block per data set (\S\ref{sec:datasets}) \\
\bottomrule
\end{tabularx}
\end{table}

\begin{note}
Changing any \code{[model]} setting means new CR tables. These are built automatically and kept in their own
folder, so the tables of the default settings are never overwritten. Check the number with \code{--plan} first.
\end{note}

% =============================================================================================
\section{Running the model}
% =============================================================================================
\begin{itemize}
  \item \code{python Main.py <input file> --plan}: check and plan only. Use it every time you change something.
  \item \code{python Main.py <input file>}: full run. It builds the missing libraries and tables, fits, and writes
        all outputs.
  \item \code{python Main.py} with no argument runs the input file named in \code{DEFAULT_INPUT} at the top of
        \code{Main.py}.
\end{itemize}

\textbf{Run time.} Everything expensive is cached and reused:
\begin{itemize}
  \item \textbf{BOLSIG+ libraries}: one per (admixture \%, pressure, $T_g$) for Ar/N$_2$ EEDFs, or per (pressure,
        $T_g$) for the pure-Ar EEDF. They take about 1--20\,min each, with 6 built in parallel, and are stored in
        \code{InputData/Bolsig/}.
  \item \textbf{CR tables}: one per condition and EEDF library, about 2--3\,min each, 6 in parallel. They are stored
        in \code{Experimental_Data/Output/AnalyticFivePFit/tables}.
  \item \textbf{Line fits} of the spectra: 1--2\,min per data set the first time. They are stored in
        \code{<output_folder>/<data set>/measurements/}, so a \emph{new} output folder fits the lines again.
  \item With everything cached, a data set takes a few minutes.
\end{itemize}
The settings of every run are written to \code{<output_folder>/settings_used.json}.

% =============================================================================================
\section{Choosing the physics}
% =============================================================================================
\subsection{The EEDF}\label{sec:eedf}
The EEDF is a BOLSIG+ solution in the 2.42\,GHz field, on an E/N grid of 3--1000\,Td. The fit picks the E/N.
\begin{itemize}
  \item \code{"pure_Ar"}: the pure-Ar EEDF. The admixture only acts in the CR model (quenching, dilution). This
        is the default.
  \item \code{"ArN2_vib_saturated"}: Ar/N$_2$ EEDF with the N$_2$ vibrations saturated (no net energy loss to
        them).
  \item \code{"ArN2_vib_Tv"}: Ar/N$_2$ EEDF with N$_2$(v) in a Boltzmann distribution at \code{Tv_K}. With
        \code{Tv_K = "LAS"}, $T_v$ is chosen per condition from the LAS 1s5, scanning \code{Tv_grid_K}.
  \item \textbf{Your own}: run \code{python CreateEEDFScheme.py} and answer the questions (gases, cross-section
        files, reactions to leave out, vibrations, field, E/N values, LAS or fixed $T_g$, a name and a comment).
        Then use \code{model = "<name>"}. All schemes, with their comments, are listed in
        \code{InputData/EEDFSchemes/EEDF_schemes.md}.
\end{itemize}

\subsection{The electron density}
\code{[fit] Ne_pin} can take four values:
\begin{itemize}
  \item \code{"free"}: $N_e$ is fitted together with E/N (2D posterior).
  \item \code{"free_else_nc"}: free, but pinned at the critical density $n_c = 7.26\times10^{16}$\,m$^{-3}$ where
        the free fit runs to the top of the $N_e$ grid (this happens at 0--1\,\% N$_2$).
  \item \code{"nc"}: pinned at $n_c$.
  \item a number, e.g. \code{1.5e17}: pinned at that value.
\end{itemize}
\code{fits.csv} always also has the free-fit $N_e$ (\code{Ne_free}) and the cost of pinning at $n_c$
(\code{dchi2_nc}). \code{nc_penalty = true} adds a soft prior against $N_e$ far above $n_c$ (off by default).

\subsection{The admixed gas}
\code{[model] admixture_gas} (default \code{"N2"}) sets which gas quenches the Ar levels in the CR model. Only the
quenching rates change; dilution, trapping and diffusion are the same for every gas. The rates are read from the
two quenching tables:
\begin{itemize}
  \item all four 4s levels: N$_2$, H$_2$, O$_2$;
  \item 4s metastables only (a 1s5/1s3 average): NO, CO$_2$, CF$_4$, CHF$_3$, CH$_4$, C$_2$H$_2$, C$_2$H$_4$,
        C$_2$H$_6$, SF$_6$, NO$_2$, F$_2$, Cl$_2$.
\end{itemize}
The run \textbf{stops} if the gas has no metastable rate (e.g.\ He, Ne, Kr, Xe, CO). If only the resonant 4s rates
are missing, it checks whether they matter, by comparing their estimated quenching with the trapped radiative
decay. Below 5\,\% it only warns; above 5\,\% it stops, which happens at our 1\,Torr conditions. The EEDF is
chosen separately, and a warning is printed when the EEDF's admixture differs from \code{admixture_gas}.

\begin{note}
There is deliberately no ``response tilt'' option: the intensity calibration is trusted.
\end{note}

% =============================================================================================
\section{The results}\label{sec:results}
% =============================================================================================
\subsection{Files}
\begin{table}[h]
\centering\small
\begin{tabularx}{\linewidth}{@{}>{\raggedright\arraybackslash}p{6.6cm}L@{}}
\toprule
File & Contents \\
\midrule
\multicolumn{2}{@{}l}{\emph{in} \code{<output_folder>/}} \\
\code{fits_all.csv} & one row per condition, all data sets: E/N, $T_{e,\mathrm{eff}}$, $N_e$ with bands, $\chi^2$/dof, CR 1s densities, LAS 1s5 \\
\code{actinometry_per_condition_all.csv} & $n_N/n_{Ar}$, N-atom density, dissociation per condition and N\,I line \\
\code{results.pkl} & the results dict \code{RES} (below) \\
\code{overview_*.png} & all data sets of one sweep type together \\
\code{<LAS table>_with_OES.csv} & the LAS table with the OES results added (densities in cm$^{-3}$) \\
\code{settings_used.json} & exactly which settings the run used \\
\midrule
\multicolumn{2}{@{}l}{\emph{in} \code{<output_folder>/<data set>/}} \\
\code{fits.csv}, \code{residuals.csv} & the fits and the line residuals of the data set \\
\code{actinometry_per_condition.csv}, \code{actinometry_per_pair.csv}, \code{line_ratios.csv} & actinometry per condition and per N\,I/Ar\,I pair; measured line ratios \\
\code{figures/} & metastables vs LAS, E/N and $T_e$, $N_e$, EEDFs, residuals, actinometry, raw-line zooms \\
\code{measurements/} & the cached line fits of the spectra \\
\bottomrule
\end{tabularx}
\end{table}

Main columns of \code{fits.csv}: \code{x} (sweep value), \code{N2}, \code{p_Torr}, \code{P_W}, \code{Tg};
\code{EN_nc}, \code{Te_nc} (with \code{_lo}/\code{_hi}); \code{Ne_fit} (with \code{_lo}/\code{_hi}) and
\code{Ne_mode}; \code{chi2_red_nc}; \code{n_1s5}, \code{n_1s4}, \code{n_1s3}, \code{n_1s2} [m$^{-3}$];
\code{las_1s5} and \code{ratio_cr_las}; \code{Ne_free}, \code{dchi2_nc}. The suffix \code{_nc} means ``the fit
of the chosen $N_e$ mode'', which is pinned or free.

\subsection{The results dict \texttt{RES}}
\code{Main.py} returns \code{RES}, and it is saved as \code{<output_folder>/results.pkl}. It is ordered by data
set and condition:
\begin{lstlisting}
RES["80W_1Torr"]["0.33 % N2"] = {
    "condition":          data set, sweep, x, pressure, power, admixture gas and %, Tg
    "files":              the spectra of this condition
    "fit":                the fits.csv row (E/N, Te_eff, Ne, chi2, CR 1s densities, LAS 1s5, ...)
    "scale_per_spectrum": {file: s}       (ln I_measured = ln I_model + s)
    "line_residuals":     table of every measured line of every spectrum vs the model
    "eedf":               the EEDF at the fitted E/N (E_eV, EEPF, Te_eff_eV, library)
    "CR":                 {level: ...}  the CR model solved at the fit: for every state its
                          density_m^-3, Loss_s^-1, Production_m^-3s^-1, kind, energy, rates,
                          GasQuenching (what was read from the quenching tables), ...
    "CR_solver":          Ne and E/N of the solve, convergence, the CR table used
    "actinometry":        actinometry rows of this condition (N2 conditions)
}
RES["by_file"]["Percent_0_33__2"]   ->  the same entry, looked up by spectrum file
\end{lstlisting}
The condition keys are the sweep value with its unit: \texttt{"0.33 \% N2"}, \code{"850 mTorr"},
\code{"35 W"}. Reading it later, without running anything:
\begin{lstlisting}[language=Python]
import pandas as pd
RES = pd.read_pickle(r"Experimental_Data/Output/MyFirstRun/results.pkl")
e = RES["80W_1Torr"]["0.33 % N2"]
print(e["fit"]["EN_nc"], e["fit"]["Ne_fit"], e["fit"]["chi2_red_nc"])
dens = {lvl: s["density_m^-3"] for lvl, s in e["CR"].items()}       # all states
print(dens["4s1"], dens["4s3"])                                      # 1s5, 1s3 metastables
print(RES["by_file"]["Percent_0_33__2"]["condition"])
\end{lstlisting}
The densities in \code{CR} come from a new CR solve at the fitted EEDF and $N_e$. The densities in \code{fit} are
interpolated in the CR table, and the two agree to within about 1\,\%.

% =============================================================================================
\section{Troubleshooting}
% =============================================================================================
\begin{table}[h]
\centering\small
\begin{tabularx}{\linewidth}{@{}>{\raggedright\arraybackslash}p{5.6cm}L@{}}
\toprule
Message / problem & What to do \\
\midrule
\emph{unknown section} / \emph{has no setting} & a typo in the input file; the message lists the allowed names \\
\emph{no calibrated spectrum matches the file pattern} & check \code{folder}, \code{files} and \code{[spectra] export_flags}; export the spectra with the intensity calibration \\
\emph{N spectra without the intensity calibration left out} & these exports are raw; they are skipped \\
\code{--plan} wants many libraries / tables & normal the first time for new pressures or $T_g$; a coarser \code{Tg_step_K} needs fewer \\
\emph{BOLSIG+ libraries failed} / \emph{CR tables failed} & logs in \code{Experimental_Data/Output/LASTgComparison/logs} and \code{.../AnalyticFivePFit/logs} \\
\emph{no quenching rate ... by admixture\_gas} & the gas is missing from the quenching tables: add its rates, or check the name \\
$N_e$ at $3\times10^{19}$\,m$^{-3}$ (grid edge) & the free fit can run to the dense branch at low N$_2$: use \code{Ne_pin = "free_else_nc"} or \code{"nc"} \\
results from an earlier run left over & \code{[run] overwrite = true} clears a data set's old results before its run; delete \code{<set>/measurements/} to refit the lines \\
figures do not open & set \code{[run] show_plots = true} (the run then waits until you close them) \\
\bottomrule
\end{tabularx}
\end{table}

% =============================================================================================
\section{Where things are in the code}
% =============================================================================================
You do not need this to run the model. Every file in \code{pipeline/} starts with a header saying what it takes
in, what it gives out and which files use it, and \code{pipeline/README.md} has a map of the files.
\begin{itemize}
  \item \code{Main.py}: reads and checks the input file, and passes the settings on (\code{DEFAULTS} lists them all).
  \item \code{pipeline/RunDatasets.py}: the analysis of each data set and the results dict.
  \item \code{pipeline/CRModelSolver.py}: the CR model (balance equations of the 51 Ar levels).
  \item \code{pipeline/CRTableFit.py}: CR tables over (E/N, $N_e$), the line fit, and admixture quenching.
  \item \code{pipeline/LASAndEEDFLibraries.py}: LAS $T_g$ and 1s5, and the BOLSIG+ libraries.
  \item \code{pipeline/PhysicsHelpers.py}: atomic data, EEDFs, BOLSIG+, radiation trapping.
\end{itemize}
To add a new option: add it to \code{DEFAULTS} in \code{Main.py} and to \code{input_example.toml} with a comment,
then check it with \code{--plan}. The old module names (\code{CRFitNeTe}, \code{HelperFunctions},
\code{N2SweepPureArEEDF}, \dots) still work in older scripts.

\end{document}
